\begin{abstract}
We characterize minimax squared-error precision for the population all-treated-versus-all-control contrast under heterogeneous additive interference and sampled network information. The experiment assigns treatment independently with probability one half, retains each true directed off-diagonal arrow independently with known probability \(q\), the edge-retention probability, and observes the retained graph, every assignment, and every noiseless realized outcome. Fixed schedules have population size \(n\), common in- and out-degree bound \(d\), and absolute coefficient mass at most one for each outcome, including its baseline and own-treatment coefficients. Uniformly over integers \(n\ge4\) and \(1\le d\le n-1\), and retention probabilities \(0<q\le1\), minimax squared risk is of order
\[
\min\left\{1,\frac{d^2}{n[1-(1-q)^d]}\right\}
\asymp
\min\left\{1,\frac{d^2}{n}+\frac{d}{nq}\right\},
\]
with universal comparison constants; at zero retention it is of constant order. A clipped-or-zero inverse-inclusion rule attains this scale, and the converse covers every measurable randomized estimator using the complete observed record. For degree sequences \(d_n\) and retention sequences \(q_n\), uniform consistency is possible exactly when \(d_n^2/n\to0\) and \(d_n/(nq_n)\to0\), interpreting the second ratio as infinite at zero retention. Retention of order \(1/d\) or greater attains the supplied-graph precision order. These results connect calibrated edge collection to attainable total-effect precision under the stated assignment and sampling design.
\end{abstract}

\section{Introduction}

This paper characterizes the optimal squared-error precision for estimating a population total treatment effect from one randomized experiment with a sampled interference graph. The target compares assigning treatment to everyone with assigning treatment to no one, averaging the resulting outcome differences over a fixed population. We study heterogeneous additive responses under bounded absolute coefficient mass and bounded directed in- and out-degrees. Treatment is assigned independently with probability one half, and each true off-diagonal interference arrow is independently retained with a known common probability. Under this experiment, we establish matching minimax risk bounds, construct an observable attaining rule, and characterize exactly when uniformly consistent estimation is possible.

The precision question arises when network information is collected alongside an experiment. Assignments and realized outcomes are available for the full population, while recorded arrows reveal a sampled portion of the interference structure. The retained graph contributes information about which treatments enter each outcome; the outcomes themselves can also inform those associations. An optimal-precision analysis therefore evaluates the joint information in graph labels, assignments, and outcomes. Our minimax criterion allows every measurable randomized estimator to use that complete record, as specified in \cref{obj:synth_3,obj:def:risk}.

The schedule restrictions make the causal target and information comparison precise. Write \(n\), the population size, and \(d\), the common off-diagonal in-degree and out-degree bound, with \(n\ge4\) and \(1\le d\le n-1\). Each outcome is the sum of a fixed baseline, an own-treatment effect, and source-specific additive spillovers. For each recipient, the sum of the absolute baseline, own-treatment, and spillover coefficients is at most one. Coefficients may be heterogeneous and have either sign. The graph and coefficients remain fixed during treatment randomization and edge collection. This is the schedule class in \cref{obj:def:model}. Its correspondence with the published interaction-order-one SNIPE model adjoins every own-treatment coordinate and translates the off-diagonal degree bound into an augmented degree bound of \(d+1\), as recorded in \cref{obj:lem:published-model}.

Let \(q\), the known edge-retention probability, range over \([0,1]\). The principal result in \cref{obj:thm:observed-record-precision-frontier} establishes, uniformly over the stated parameter domain, that the minimax squared risk \(R_{n,d}(q)\) satisfies
\[
R_{n,d}(q)\asymp
\begin{cases}
\displaystyle
\min\left\{1,\frac{d^2}{n[1-(1-q)^d]}\right\},&q>0,\\[6pt]
1,&q=0,
\end{cases}
\]
with universal comparison constants. For positive retention, the equivalent expression
\[
R_{n,d}(q)\asymp
\min\left\{1,\frac{d^2}{n}+\frac{d}{nq}\right\}
\]
separates a degree contribution from an edge-collection contribution. The cap at one describes saturation of worst-case squared risk under the bounded target range.

The upper comparison is constructive. The inverse-inclusion score in \cref{obj:def:audit-score} combines each realized outcome with its own-treatment sign and the treatment signs of retained sources, weighting recorded arrows by \(q^{-1}\). For every fixed admissible schedule at positive retention, this score is unbiased over assignment and audit randomness. Its variance decomposes into assignment variation with the graph supplied and an additional audit term, as shown in \cref{obj:lem:score-audit}. The attaining rule clips the score to the target interval when its risk envelope is below one and uses zero otherwise. \Cref{obj:thm:observable-upper,obj:thm:observed-record-precision-frontier} establish its uniform guarantee and minimax order.

The lower comparison evaluates the same complete observation experiment. It uses opposite-sign priors over fixed schedules, with baseline and source-allocation draws occurring before the experiment. The resulting likelihood preserves retained-edge labels, all treatment coordinates, repeated recipient outcomes, and the treatment-count constraint on hidden source allocations. Translation bounds, all-order Walsh energy control, and hidden-allocation contraction yield the testing comparison in \cref{obj:thm:block-testing-scale}. Together with the two-prior risk reduction in \cref{obj:lem:full-record-two-prior-risk} and the supplied-graph lower comparison in \cref{obj:lem:supplied-block-record-lower}, these arguments establish the converse against arbitrary measurable randomized uses of the record throughout the admissible degree and retention ranges.

The characterization gives two collection implications. First, retention of order \(1/d\) or greater attains the supplied-graph risk order \(\min\{1,d^2/n\}\); below this transition, the collection contribution determines the uncapped scale. Second, along admissible population sequences with degree bounds \(d_n\) and retention probabilities \(q_n\), uniform consistency is possible exactly when
\[
d_n=o(\sqrt n),
\qquad
q_n>0\ \text{eventually},
\qquad
\frac{nq_n}{d_n}\longrightarrow\infty.
\]
Thus consistency can coexist with vanishing retention, and retention of order \(1/d\) or greater suffices to match supplied-graph precision. These conclusions follow from \cref{obj:thm:observed-record-precision-frontier,obj:thm:frontier-answer}.

The closest estimator comparison is known-neighborhood SNIPE, which studies the same population contrast under Bernoulli assignment and bounded-order polynomial responses \citep[Sections 3--5; Theorems 1--2]{CortezRodriguezEichhornYu2023SNIPE}. Its published variance bound specializes to order \((d+1)^3/n\) under the additive, half-Bernoulli, unit-envelope conventions used here; its published minimax construction has its own treatment-probability conditions. The sharper supplied-graph degree benchmark of \citet[Theorem \texttt{thm:degree-frontier}; non-archival research note]{CausalSmithKnownGraph2026} is restated for the present conventions in \cref{obj:thm:supported-boundaries}. Design-based minimax analysis with a supplied network further studies joint optimization over assignments and estimators under a specified response envelope \citep[Corollary 5.4]{KandirosHarshawSavje2026DesignBasedMinimax}. Our new contribution fixes the assignment law and characterizes sampled-network precision, including a complete-record converse and the joint degree--retention conditions for consistency.

Known network-error probabilities also support correction of exposure-specific means \citep[Section 3; Theorems 4--7]{LiSussmanKolaczyk2021NetworkNoise}, while supplied neighborhood supersets support unbiased total-effect estimation under additive and polynomial responses \citep[Theorems 1, 4, and 8; Appendix B.1]{ShankarEtAl2024UNITE}. These approaches clarify the role of network information in causal estimation. The present result characterizes worst-case total-effect precision under independent retention of true arrows and heterogeneous additive coefficient restrictions. Under bounded outcomes and first- and second-order influence restrictions, the pseudoinverse estimator of \citet[Assumptions 1--2 and Theorems 1--2]{JiangDengWangWang2024ApproximateNetworks} has asymptotic variance \(O\{d_{\mathcal G}^2/[n\pi(1-\pi)]\}\), with the same estimator-specific order for constant outcomes on a regular surrogate graph. Our result instead optimizes over all estimators under the calibrated retention experiment.

Proofs of the results appear in the appendices.

The paper proceeds through related work, the population model and observation design, minimax precision and attainable estimation, and the edge-collection comparison with the supplied-graph benchmark. A limitations and future-work discussion follows. The appendices develop the score calculations, complete-record block experiments, analytic and testing bounds, and minimax arguments, followed by the verification note and proofs of the main results.
